\documentclass[12pt]{article}
\usepackage[a4paper,margin=1in]{geometry}
\usepackage[utf8]{inputenc}
\usepackage[T1]{fontenc}
\usepackage{amsmath,amssymb}
\usepackage{textcomp}
\usepackage{parskip}
\begin{document}
\begin{center}
{\Large \textbf{2020 STEP III}}
\end{center}

\subsection*{Page 1}
\begin{flushleft}
\bigskip\textbf{Section A: Pure Mathematics}\par
\bigskip\textbf{Question 1.} For non-negative integers a and b, let\par
\[Z \pi\]
\[I(a, b) = 2 cosa x cos bx dx.\]
\textbf{(i)} Show that for positive integers a and b,\par
a\par
\[I(a, b) = I(a - 1, b - 1).\]
a + b\par
\textbf{(ii)} Prove by induction on n that for non-negative integers n and m,\par
\[Z \pi 2 cosn x cos(n + 2m + 1)x dx = (-1)m 2nn!(2m)!(n + m)! .\]
m!(2n + 2m + 1)!\par
\bigskip\textbf{Question 2.} The curve C has equation sinh x + sinh y = 2k, where k is a positive constant.\par
\[d2y\]
\textbf{(i)} Show that the curve C has no stationary points and that = 0 at the point (x, y)\par
\[dx2\]
on the curve if and only if\par
\[1 + sinh x sinh y = 0.\]
Find the co-ordinates of the points of inflection on the curve C, leaving your answers\par
in terms of inverse hyperbolic functions.\par
\textbf{(ii)} Show that if (x, y) lies on the curve C and on the line x + y = a, then\par
\[e2x(1 - e-a) - 4kex + (ea - 1) = 0\]
\[and deduce that 1 < cosh a 6 2k2 + 1.\]
\textbf{(iii)} Sketch the curve C.\par
\end{flushleft}
\subsection*{Page 2}
\begin{flushleft}
\bigskip\textbf{Question 3.} Given distinct points A and B in the complex plane, the point G is defined to be the\par
AB\par
centroid of the triangle ABK, where the point K is the image of B under rotation about A\par
\[through a clockwise angle of 1 \pi.\]
Note: if the points P , Q and R are represented in the complex plane by p, q and r, the\par
centroid of triangle P QR is defined to be the point represented by 1 (p + q + r).\par
\textbf{(i)} If A, B and G are represented in the complex plane by a, b and g , show that\par
AB ab\par
\[g = \sqrt (ωa + ω∗b),\]
ab\par
\[i\pi\]
\[where ω = e 6 .\]
\textbf{(ii)} The quadrilateral Q has vertices A, B, C and D, in that order, and the quadrilateral\par
Q has vertices G , G , G and G , in that order. Using the result in part (i),\par
2 AB BC CD DA\par
show that Q is a parallelogram if and only if Q is a parallelogram.\par
1 2\par
\textbf{(iii)} The triangle T has vertices A, B and C and the triangle T has vertices G , G\par
1 2 AB BC\par
and G . Using the result in part (i), show that T is always an equilateral triangle.\par
CA 2\par
\end{flushleft}
\subsection*{Page 3}
\begin{flushleft}
\bigskip\textbf{Question 4.} The plane Π has equation r . n = 0 where n is a unit vector. Let P be a point with position\par
vector x which does not lie on the plane Π. Show that the point Q with position vector\par
x - (x . n)n lies on Π and that P Q is perpendicular to Π.\par
\textbf{(i)} Let transformation T be a reflection in the plane ax+by+cz = 0, where a2 +b2 +c2 = 1.\par
\[ 1   b^2 + c^2 - a^2 \]
\[Show that the image of i =  0  under T is  -2ab , and find the images\]
0 -2ac\par
of j and k under T .\par
Write down the matrix M which represents transformation T .\par
\textbf{(ii)} The matrix\par
 \par
0.64 0.48 0.6\par
0.48 0.36 -0.8\par
0.6 -0.8 0\par
represents a reflection in a plane. Find the cartesian equation of the plane.\par
\textbf{(iii)} The matrix N represents a rotation through angle \pi about the line through the origin\par
 \par
a\par
\[parallel to  b , where a^2 + b^2 + c^2 = 1. Find the matrix N.\]
\textbf{(iv)} Identify the single transformation which is represented by the matrix NM.\par
\end{flushleft}
\subsection*{Page 4}
\begin{flushleft}
n\par
X\par
\bigskip\textbf{Question 5.} Show that for positive integer n, xn - yn = (x - y) xn-ryr-1.\par
\[r=1\]
\textbf{(i)} Let F be defined by\par
\[F(x) = for x 6= 0, k\]
xn(x - k)\par
\[where n is a positive integer and k 6= 0.\]
(a) Given that\par
A f(x)\par
\[F(x) = + ,\]
x - k xn\par
where A is a constant and f(x) is a polynomial, show that\par
1 \gamma \gamma x \deltan\delta\par
\[f(x) = 1 - .\]
x - k k\par
Deduce that\par
n\par
1 1 X 1\par
\[F(x) = - .\]
kn(x - k) k kn-rxr\par
\[r=1\]
(b) By differentiating xn F(x), prove that\par
n\par
1 1 n X n - r\par
\[= - + .\]
xn(x - k)2 kn(x - k)2 xkn(x - k) kn+1-rxr+1\par
\[r=1\]
\textbf{(ii)} Hence evaluate the limit of\par
Z N 1\par
dx\par
\[x^3(x - 1)2\]
as N \to \infty, justifying your answer.\par
\end{flushleft}
\subsection*{Page 5}
\begin{flushleft}
\[\sqrt\]
\bigskip\textbf{Question 6.} (i) Sketch the curve y = cos x + cos 2x for - 1 \pi 6 x 6 1 \pi.\par
4 4\par
\textbf{(ii)} The equation of curve C in polar co-ordinates is\par
\[\sqrt\]
\[r = cos θ + cos 2θ - 1 \pi 6 θ 6 1 \pi.\]
4 4\par
Sketch the curve C .\par
\textbf{(iii)} The equation of curve C in polar co-ordinates is\par
\[r^2 - 2r cos θ + sin2 θ = 0 - 1 \pi 6 θ 6 1 \pi.\]
4 4\par
\[Find the value of r when θ = ± 1 \pi.\]
Show that, when r is small, r \approx 1 θ2.\par
\[Sketch the curve C , indicating clearly the behaviour of the curve near r = 0 and\]
\[near θ = ± 1 \pi.\]
\[\pi\]
\[Show that the area enclosed by curve C and above the line θ = 0 is \sqrt .\]
2 2\par
\bigskip\textbf{Question 7.} (i) Given that the variables x, y and u are connected by the differential equations\par
du dy\par
\[+ f(x)u = h(x) and + g(x)y = u,\]
dx dx\par
show that\par
\[d2y dy\]
\[+ (g(x) + f(x)) + (g^0(x) + f(x)g(x))y = h(x). (1)\]
\[dx2 dx\]
\textbf{(ii)} Given that the differential equation\par
\[d2y \lambda 4 \mu dy \lambda 2 2 \mu\]
\[+ 1 + + + y = 4x + 12 (2)\]
\[dx2 x dx x x^2\]
can be written in the same form as (1), find a first order differential equation which is\par
satisfied by g(x).\par
\[If g(x) = kxn, find a possible value of n and the corresponding value of k.\]
dy\par
\[Hence find a solution of (2) with y = 5 and = -3 at x = 1.\]
dx\par
\end{flushleft}
\subsection*{Page 6}
\begin{flushleft}
\bigskip\textbf{Question 8.} A sequence u , for integer k > 1, is defined as follows.\par
k\par
\[u = 1\]
\[u = u for k > 1\]
2k k\par
\[u = u + u for k > 1\]
2k+1 k k+1\par
\textbf{(i)} Show that, for every pair of consecutive terms of this sequence, except the first pair,\par
the term with odd subscript is larger than the term with even subscript.\par
\textbf{(ii)} Suppose that two consecutive terms in this sequence have a common factor greater\par
than one. Show that there are then two consecutive terms earlier in the sequence\par
which have the same common factor. Deduce that any two consecutive terms in this\par
sequence are co-prime (do not have a common factor greater than one).\par
\textbf{(iii)} Prove that it is not possible for two positive integers to appear consecutively in the\par
same order in two different places in the sequence.\par
\textbf{(iv)} Suppose that a and b are two co-prime positive integers which do not occur consec-\par
\[utively in the sequence with b following a. If a > b, show that a - b and b are two\]
co-prime positive integers which do not occur consecutively in the sequence with b\par
\[following a - b, and whose sum is smaller than a + b. Find a similar result for a < b.\]
\textbf{(v)} For each integer n > 1, define the function f from the positive integers to the positive\par
u\par
n\par
\[rational numbers by f(n) = . Show that the range of f is all the positive rational\]
u\par
n+1\par
numbers, and that f has an inverse.\par
\end{flushleft}
\subsection*{Page 7}
\begin{flushleft}
\bigskip\textbf{Section B: Mechanics}\par
\bigskip\textbf{Question 9.} Two inclined planes Π and Π meet in a horizontal line at the lowest points of both planes\par
1 2\par
and lie on either side of this line. Π and Π make angles of α and β, respectively, to the\par
1 2\par
\[horizontal, where 0 < β < α < 1 \pi.\]
A uniform rigid rod P Q of mass m rests with P lying on Π and Q lying on Π so that the\par
1 2\par
rod lies in a vertical plane perpendicular to Π and Π with P higher than Q.\par
1 2\par
\textbf{(i)} It is given that both planes are smooth and that the rod makes an angle θ with the\par
\[horizontal. Show that 2 tan θ = cot β - cot α.\]
\textbf{(ii)} It is given instead that Π is smooth, that Π is rough with coefficient of friction µ and\par
1 2\par
that the rod makes an angle φ with the horizontal. Given that the rod is in limiting\par
equilibrium, with P about to slip down the plane Π , show that\par
µ\par
\[tan θ - tan φ =\]
(µ + tan β) sin 2β\par
\[where θ is the angle satisfying 2 tan θ = cot β - cot α.\]
\bigskip\textbf{Question 10.} A light elastic spring AB, of natural length a and modulus of elasticity kmg, hangs vertically\par
with one end A attached to a fixed point. A particle of mass m is attached to the other\par
\[end B. The particle is held at rest so that AB > a and is released.\]
Find the equation of motion of the particle and deduce that the particle oscillates vertically.\par
\[2\pi\]
\[If the period of oscillation is , show that kg = aΩ2.\]
Ω\par
Suppose instead that the particle, still attached to B, lies on a horizontal platform which\par
\[2\pi\]
performs simple harmonic motion vertically with amplitude b and period .\par
ω\par
At the lowest point of its oscillation, the platform is a distance h below A.\par
Let x be the distance of the particle above the lowest point of the oscillation of the platform.\par
When the particle is in contact with the platform, show that the upward force on the particle\par
from the platform is\par
mg + mΩ2(a + x - h) + mω2(b - x).\par
Given that ω < Ω, show that, if the particle remains in contact with the platform throughout\par
its motion,\par
\lambda 1 \mu ω2b\par
h 6 a 1 + + .\par
k Ω2\par
\[Find the corresponding inequality if ω > Ω.\]
Hence show that, if the particle remains in contact with the platform throughout its motion,\par
it is necessary that\par
\lambda \mu\par
h 6 a 1 + + b,\par
k\par
whatever the value of ω.\par
\end{flushleft}
\subsection*{Page 8}
\begin{flushleft}
\bigskip\textbf{Section C: Probability and Statistics}\par
\bigskip\textbf{Question 11.} The continuous random variable X is uniformly distributed on [a, b] where 0 < a < b.\par
\textbf{(i)} Let f be a function defined for all x \in [a, b]\par
\[• with f(a) = b and f(b) = a,\]
• which is strictly decreasing on [a, b],\par
\[• for which f(x) = f-1(x) for all x \in [a, b].\]
\[The random variable Y is defined by Y = f(X). Show that\]
b - f(y)\par
\[P(Y 6 y) = for y \in [a, b].\]
b - a\par
Find the probability density function for Y and hence show that\par
Z b 2xf(x)\par
\[E(Y 2) = -ab + dx.\]
b - a\par
a\par
1 1 1 1 1 1\par
\textbf{(ii)} The random variable Z is defined by + = where = + . By finding the\par
Z X c c a b\par
variance of Z, show that\par
\lambda \mu\par
b - c b - a\par
\[ln < .\]
a - c c\par
\bigskip\textbf{Question 12.} A and B both toss the same biased coin. The probability that the coin shows heads is p,\par
\[where 0 < p < 1, and the probability that it shows tails is q = 1 - p.\]
Let X be the number of times A tosses the coin until it shows heads. Let Y be the number\par
of times B tosses the coin until it shows heads.\par
\textbf{(i)} The random variable S is defined by S = X + Y and the random variable T is the\par
\[maximum of X and Y . Find an expression for P(S = s) and show that\]
\[P(T = t) = pqt-1(2 - qt-1 - qt).\]
\textbf{(ii)} The random variable U is defined by U = |X - Y |, and the random variable W is the\par
\[minimum of X and Y . Find expressions for P(U = u) and P(W = w).\]
\textbf{(iii)} Show that P(S = 2 and T = 3) 6= P(S = 2) \times P(T = 3).\par
\textbf{(iv)} Show that U and W are independent, and show that no other pair of the four variables\par
S, T , U and W are independent.\par
\end{flushleft}
\subsection*{Page 9}
\begin{flushleft}
BLANK PAGE\par
\end{flushleft}
\end{document}
